% Prepared from templates/tutorial.tex; Tutorial 6, 2026-11-06 (Week 12).
% Source: Tutorial 6 - Default Logic and Fuzzy Logic.pdf, printed/PDF p. 1 of 1.
% Original questions 5.1--5.3; formulas and graph checked visually.
% No official solution supplied. Figures are redrawn from source breakpoints.
\begingroup
\renewcommand{\recordid}[1]{}
\renewcommand{\relatedtopic}[2]{\hyperref[topic:#1]{#2}}
\exercisegroup{SC3000-TUT06}{Tutorial 06：Default Logic 与 Fuzzy Logic}
\label{tutorial:SC3000-TUT06}

\noteintro{\textbf{日期与资料：}2026-11-06（Week 12）；Tutorial 6 -- Default Logic and Fuzzy Logic，p.~1（全文），原题号 5.1--5.3。按原题公式与图形整理推导；未提供官方参考答案。}

\linespread{1.1}\selectfont
\setlength{\parskip}{0.35em}

\exercisequestion{SC3000-TUT06-Q01}{5.1：求默认理论的全部 extensions}
\textbf{对应讲义：}\relatedtopic{SC3000-L10-T01}{默认规则}、\relatedtopic{SC3000-L10-T02}{Extension 与猜测验证}。

\textbf{题目。}求以下两个 default theories（默认理论）$T=\langle\Delta,\Phi\rangle$ 的全部 extensions（扩展）：
\begin{align*}
\text{(1)}\quad
\Delta&=\left\{\frac{R(x):\neg P(x)}{\neg P(x)},\frac{Q(x):P(x)}{P(x)}\right\},
&\Phi&=\{R(N)\land Q(N)\};\\
\text{(2)}\quad
\Delta&=\left\{\frac{Summer:\neg Rain}{Sun\_Shining}\right\},
&\Phi&=\{\neg Sun\_Shining\land Summer\}.
\end{align*}

\begin{checkedsolution}
\textbf{判断方法。}规则 $\frac{\alpha:\beta}{\gamma}$ 的先决条件 $\alpha$ 须由当前构造推出；其 justification（合理性条件）对最终候选 $E$ 检查，要求 $\neg\beta\notin E$。只加入结论 $\gamma$，再取经典逻辑闭包 $Cn$。最终构造必须与候选 $E$ 相同。

\textbf{(1) 两个扩展。}以 $x=N$ 实例化，两条规则的先决条件均由 $\Phi$ 提供。得到
\[
 E_1=Cn\{R(N),Q(N),\neg P(N)\},\qquad
 E_2=Cn\{R(N),Q(N),P(N)\}.
\]
在 $E_1$ 中第一条规则可用，第二条被 $\neg P(N)$ 阻止；在 $E_2$ 中第二条可用，第一条被 $P(N)$ 阻止。若两种结论都不含，两条规则均可用，候选尚未封闭；若两种都含，经典闭包不一致，两条规则又均被阻止，从 $\Phi$ 无法构造出该候选。因此只有上述两个扩展。不同尝试顺序可以找到不同分支，求全部扩展时须保留两者。

\textbf{(2) 无扩展。}记 $S=Sun\_Shining$。按候选 $E$ 是否包含 $Rain$ 分两种情况：
\begin{itemize}
 \item 若 $Rain\notin E$，规则可用，从已知的 $Summer,\neg S$ 推出 $S$。经典逻辑闭包由矛盾推出所有公式，也包含 $Rain$，构造结果与候选不一致。
 \item 若 $Rain\in E$，规则被阻止，构造只能得到 $Cn\{Summer,\neg S\}$，其中推不出 $Rain$，仍与候选不一致。
\end{itemize}
两个分支均无法形成固定点，故不存在 extension。这里 justification 是 $\neg Rain$，结论是 $S$；已有 $\neg S$ 并不直接阻止规则，也不能据此把 $Rain$ 补成事实。Justification 用于相容性检查，不会自动加入结论集合。
\end{checkedsolution}

\exercisequestion{SC3000-TUT06-Q02}{5.2：从图形读取梯形参数与划分类型}
\textbf{对应讲义：}\relatedtopic{SC3000-L11-T01}{梯形隶属函数}、\relatedtopic{SC3000-L11-T02}{模糊划分}。

\textbf{题目。}变量 $X\in[0,21]$ 的四个 fuzzy labels（模糊标签）如下。用 $\operatorname{tpmf}[a,b,c,d]$ 描述各隶属函数，并判断它们在该论域上的划分类型。
\begin{center}
% Original redraw of Tutorial 6, p. 1, Figure Q5.2; numeric x scale.
\begin{tikzpicture}[x=0.61cm,y=1.35cm,>=Latex,font=\small]
 \draw[->] (0,0) -- (22,0) node[right] {$x$};
 \draw[->] (0,0) -- (0,1.45) node[above] {$\mu(x)$};
 \draw[gray!45,dotted] (0,1) node[left,black] {$1$} -- (21,1);
 \draw[noteBlue,thick] (0,0) -- (3,1) -- (7,0);
 \draw[verifyOrange,thick,dashed] (4,0) -- (6,1) -- (9,1) -- (12,0);
 \draw[teal!70!black,thick,dash dot] (9,0) -- (12,1) -- (16,1) -- (21,0);
 \draw[violet,thick,densely dotted] (15,0) -- (18,1) -- (21,1);
 \foreach \x in {0,3,4,6,7,9,12,15,16,18,21}{
  \draw (\x,0.03) -- (\x,-0.03) node[below] {\scriptsize $\x$};
 }
 \node[above,noteBlue] at (3,1) {$A_{11}$};
 \node[above,verifyOrange] at (7.5,1) {$A_{12}$};
 \node[above,teal!70!black] at (14,1) {$A_{13}$};
 \node[above,violet] at (19.5,1) {$A_{14}$};
\end{tikzpicture}

\end{center}
\textbf{图示：}按原题 Figure Q5.2 的转折点重绘，横轴按数值比例绘制。

\begin{checkedsolution}
$a,b,c,d$ 依次表示开始上升、升至 $1$、开始下降、降至 $0$ 的位置；$b=c$ 对应三角形。右肩形用论域右端的 $c=d$ 表示到端点仍保持为 $1$，不计算零分母的下降段。
\begin{center}
\begin{tabular}{@{}lll@{}}
\toprule
标签 & 参数 & 形状 \\
\midrule
$A_{11}$ & $[0,3,3,7]$ & 三角形 \\
$A_{12}$ & $[4,6,9,12]$ & 梯形 \\
$A_{13}$ & $[9,12,16,21]$ & 梯形 \\
$A_{14}$ & $[15,18,21,21]$ & 右肩形 \\
\bottomrule
\end{tabular}
\end{center}
按本课定义，pseudo fuzzy partition（伪模糊划分）要求整个论域上 $\sum_i\mu_{A_i}(x)=1$。对 $0\le x<3$，此处只有 $A_{11}$ 可能非零，故总和为 $x/3<1$。因此本题为一般 fuzzy partition（模糊划分），不满足伪模糊划分的归一化条件；$x=0$ 处四个标签均为零。
\end{checkedsolution}

\exercisequestion{SC3000-TUT06-Q03}{5.3：小费规则、输出隶属函数与 very 修饰}
\textbf{对应讲义：}\relatedtopic{SC3000-L11-T02}{模糊划分}、\relatedtopic{SC3000-L11-T03}{集合运算与语言修饰}、\relatedtopic{SC3000-L11-T04}{模糊规则系统}。

\textbf{题目。}服务质量 $s$、食物质量 $f$ 与小费 $t$ 的论域均为 $[0,10]$，各标签的 $\operatorname{tpmf}$ 参数如下。原题输出变量交替写为 t-payment 与 t-quality，下文统一用 $t$ 表示。
\begin{center}
\small
\begin{tabular}{@{}lll@{}}
\toprule
变量 & 标签 & $[a,b,c,d]$ \\
\midrule
s-quality & poor / good / excellent & $[0,0,4,5]$ / $[4,5,6,7]$ / $[6,7,10,10]$ \\
f-quality & lousy / delicious & $[0,0,2,3]$ / $[7,8,10,10]$ \\
t-quality & cheap / average / generous & $[0,2,2,3]$ / $[3,4,4,5]$ / $[4,5,5,9]$ \\
\bottomrule
\end{tabular}
\end{center}
\begin{enumerate}[label=(\roman*)]
 \item 绘制三组变量的模糊划分，并判断每组的划分类型。
 \item 采用下列四条规则，在 $s=3,f=7$ 时求小费的输出隶属函数：
 \begin{itemize}
  \item R1：service is poor $\Rightarrow$ tip is cheap。
  \item R2：service is excellent AND food is delicious $\Rightarrow$ tip is generous。
  \item R3：food is lousy $\Rightarrow$ tip is cheap。
  \item R4：service is good AND food is delicious $\Rightarrow$ tip is average。
 \end{itemize}
 \item 将 R1 改成 ``If service is very poor, then tip is very cheap''，说明服务与小费隶属函数如何改变。
\end{enumerate}

\begin{checkedsolution}
\textbf{(i) 绘图与分类。}下图由题给参数绘制；同一变量的各标签使用不同线型。
\begin{center}
% Original plots from Tutorial 6, p. 1, Q5.3(i); shoulders retain endpoint 1.
\begin{tikzpicture}[x=0.41cm,y=1.05cm,>=Latex,font=\scriptsize]
 \foreach \offset/\var/\title in {0/s/s-quality,12.7/f/f-quality,25.4/t/t-quality}{
  \begin{scope}[shift={(\offset,0)}]
   \node[font=\small] at (5,2.1) {\title};
   \draw[->] (0,0) -- (10.7,0) node[right] {$\var$};
   \draw[->] (0,0) -- (0,1.45) node[above] {$\mu$};
   \draw[gray!50,dotted] (0,1) node[left,black] {$1$} -- (10,1);
   \foreach \x in {0,5,10}{\draw (\x,0.035) -- (\x,-0.035) node[below] {$\x$};}
  \end{scope}
 }
 \begin{scope}
  \draw[noteBlue,thick] (0,1) -- (4,1) -- (5,0);
  \draw[verifyOrange,thick,dashed] (4,0) -- (5,1) -- (6,1) -- (7,0);
  \draw[teal!70!black,thick,dash dot] (6,0) -- (7,1) -- (10,1);
  \node[above,noteBlue] at (1.4,1.05) {poor};
  \node[above,verifyOrange] at (5.4,1.05) {good};
  \node[above,teal!70!black] at (9,1.05) {excellent};
  \foreach \x in {4,6,7}{\draw[gray!40,dotted] (\x,0) -- (\x,1);}
 \end{scope}
 \begin{scope}[shift={(12.7,0)}]
  \draw[noteBlue,thick] (0,1) -- (2,1) -- (3,0);
  \draw[verifyOrange,thick,dashed] (7,0) -- (8,1) -- (10,1);
  \node[above,noteBlue] at (1.5,1.05) {lousy};
  \node[above,verifyOrange] at (8.8,1.05) {delicious};
  \foreach \x in {2,3,7,8}{\draw (\x,0.035) -- (\x,-0.035) node[below] {$\x$};}
 \end{scope}
 \begin{scope}[shift={(25.4,0)}]
  \draw[noteBlue,thick] (0,0) -- (2,1) -- (3,0);
  \draw[verifyOrange,thick,dashed] (3,0) -- (4,1) -- (5,0);
  \draw[teal!70!black,thick,dash dot] (4,0) -- (5,1) -- (9,0);
  \node[above,noteBlue] at (1.4,1.05) {cheap};
  \node[above,verifyOrange] at (4,1.42) {average};
  \draw[verifyOrange,thin] (4,1.4) -- (4,1.05);
  \node[above,teal!70!black] at (8,1.05) {generous};
  \foreach \x in {2,3,4,9}{\draw (\x,0.035) -- (\x,-0.035) node[below] {$\x$};}
 \end{scope}
\end{tikzpicture}

\end{center}
\begin{itemize}
 \item \textbf{s-quality：伪模糊划分。}在 $[4,5]$ 上 $(5-s)+(s-4)=1$，在 $[6,7]$ 上 $(7-s)+(s-6)=1$；其余区间恰有一个标签取 $1$。整个 $[0,10]$ 上总和均为 $1$。
 \item \textbf{f-quality：一般模糊划分，存在未覆盖区间。}在 $[3,7]$ 上两个隶属度均为 $0$，不满足伪模糊划分条件。
 \item \textbf{t-quality：一般模糊划分。}例如 $0\le t<2$ 时，总和为 $t/2<1$；$t=2$ 时才等于 $1$。$[9,10]$ 上三个标签也均为 $0$。
\end{itemize}

\textbf{(ii) 模糊化、规则截断与聚合。}沿用讲义的 AND/min、min 截断及 max 聚合。首先代入输入：
\[
\bigl(\mu_{poor}(3),\mu_{good}(3),\mu_{excellent}(3)\bigr)=(1,0,0),
\qquad
\bigl(\mu_{lousy}(7),\mu_{delicious}(7)\bigr)=(0,0).
\]
$f=7$ 是 delicious 上升段的起点，隶属度为 $0$。各规则的 firing strengths（触发强度）为
\[
\alpha_1=1,\qquad\alpha_2=\min(0,0)=0,\qquad
\alpha_3=0,\qquad\alpha_4=\min(0,0)=0.
\]
令 $B_j$ 为第 $j$ 条规则的结论集合，则
\[
\mu_{out}(t)=\max_{j=1,\ldots,4}\min\{\alpha_j,\mu_{B_j}(t)\}
=\mu_{cheap}(t)=
\begin{cases}
 t/2,&0\le t\le2,\\
 3-t,&2<t<3,\\
 0,&3\le t\le10.
\end{cases}
\]
输出保留完整的 cheap 三角形。题目要求 resultant membership（结果隶属函数），到此即可；没有要求进一步解模糊成单个小费数值。

\textbf{(iii) very 对隶属度平方。}先计算原隶属度，再平方，输入坐标不变：
\[
\mu_{very\ poor}(s)=[\mu_{poor}(s)]^2,
\qquad
\mu_{very\ cheap}(t)=[\mu_{cheap}(t)]^2.
\]
具体得到
\[
\mu_{very\ poor}(s)=
\begin{cases}
 1,&0\le s\le4,\\
 (5-s)^2,&4<s<5,\\
 0,&5\le s\le10,
\end{cases}
\qquad
\mu_{very\ cheap}(t)=
\begin{cases}
 (t/2)^2,&0\le t\le2,\\
 (3-t)^2,&2<t<3,\\
 0,&3\le t\le10.
\end{cases}
\]
在本题输入 $s=3$ 时，R1' 的触发强度仍为 $1^2=1$；其余三条规则仍为 $0$，故修改后的聚合输出就是 $\mu_{very\ cheap}(t)$。原隶属度为 $0$ 或 $1$ 的位置保持不变，中间值降低；斜边变成曲线，不再是梯形或三角形的直线边。
\begin{center}
% Original comparison from Q5.3(iii); very squares membership, not x.
\begin{tikzpicture}[x=0.51cm,y=1.1cm,>=Latex,font=\small]
 \node at (5,1.75) {poor 与 very poor};
 \draw[->] (0,0) -- (10.8,0) node[right] {$s$};
 \draw[->] (0,0) -- (0,1.35) node[above] {$\mu$};
 \node[left] at (0,1) {$1$};
 \draw[noteBlue,thick] (0,1) -- (4,1) -- (5,0) -- (10,0);
 \draw[verifyOrange,thick,dashed] (0,1) -- (4,1);
 \draw[verifyOrange,thick,dashed,domain=4:5,samples=35] plot (\x,{(5-\x)^2});
 \draw[verifyOrange,thick,dashed] (5,0) -- (10,0);
 \foreach \x in {0,4,5,10}{\draw (\x,0.03) -- (\x,-0.03) node[below] {\scriptsize $\x$};}
 \begin{scope}[shift={(15,0)},x=1.1cm]
  \node at (2,1.75) {cheap 与 very cheap};
  \draw[->] (0,0) -- (4.8,0) node[right] {$t$};
  \draw[->] (0,0) -- (0,1.35) node[above] {$\mu$};
  \node[left] at (0,1) {$1$};
  \draw[noteBlue,thick] (0,0) -- (2,1) -- (3,0) -- (4.5,0);
  \draw[verifyOrange,thick,dashed,domain=0:2,samples=40] plot (\x,{(\x/2)^2});
  \draw[verifyOrange,thick,dashed,domain=2:3,samples=30] plot (\x,{(3-\x)^2});
  \draw[verifyOrange,thick,dashed] (3,0) -- (4.5,0);
  \foreach \x in {0,2,3}{\draw (\x,0.03) -- (\x,-0.03) node[below] {\scriptsize $\x$};}
  \node[font=\scriptsize] at (3.7,0.45) {$[3,10]:\ \mu=0$};
 \end{scope}
 \draw[noteBlue,thick] (3,-0.85) -- (4,-0.85) node[right,black] {原函数};
 \draw[verifyOrange,thick,dashed] (12,-0.85) -- (13,-0.85) node[right,black] {very：隶属度平方};
\end{tikzpicture}

\end{center}
\end{checkedsolution}
\endgroup
